% Options for packages loaded elsewhere
\PassOptionsToPackage{unicode}{hyperref}
\PassOptionsToPackage{hyphens}{url}
\documentclass[
]{article}
\usepackage{xcolor}
\usepackage[margin=1in]{geometry}
\usepackage{amsmath,amssymb}
\setcounter{secnumdepth}{-\maxdimen} % remove section numbering
\usepackage{iftex}
\ifPDFTeX
  \usepackage[T1]{fontenc}
  \usepackage[utf8]{inputenc}
  \usepackage{textcomp} % provide euro and other symbols
\else % if luatex or xetex
  \usepackage{unicode-math} % this also loads fontspec
  \defaultfontfeatures{Scale=MatchLowercase}
  \defaultfontfeatures[\rmfamily]{Ligatures=TeX,Scale=1}
\fi
\usepackage{lmodern}
\ifPDFTeX\else
  % xetex/luatex font selection
\fi
% Use upquote if available, for straight quotes in verbatim environments
\IfFileExists{upquote.sty}{\usepackage{upquote}}{}
\IfFileExists{microtype.sty}{% use microtype if available
  \usepackage[]{microtype}
  \UseMicrotypeSet[protrusion]{basicmath} % disable protrusion for tt fonts
}{}
\makeatletter
\@ifundefined{KOMAClassName}{% if non-KOMA class
  \IfFileExists{parskip.sty}{%
    \usepackage{parskip}
  }{% else
    \setlength{\parindent}{0pt}
    \setlength{\parskip}{6pt plus 2pt minus 1pt}}
}{% if KOMA class
  \KOMAoptions{parskip=half}}
\makeatother
\usepackage{color}
\usepackage{fancyvrb}
\newcommand{\VerbBar}{|}
\newcommand{\VERB}{\Verb[commandchars=\\\{\}]}
\DefineVerbatimEnvironment{Highlighting}{Verbatim}{commandchars=\\\{\}}
% Add ',fontsize=\small' for more characters per line
\newenvironment{Shaded}{}{}
\newcommand{\AlertTok}[1]{\textcolor[rgb]{1.00,0.00,0.00}{\textbf{#1}}}
\newcommand{\AnnotationTok}[1]{\textcolor[rgb]{0.38,0.63,0.69}{\textbf{\textit{#1}}}}
\newcommand{\AttributeTok}[1]{\textcolor[rgb]{0.49,0.56,0.16}{#1}}
\newcommand{\BaseNTok}[1]{\textcolor[rgb]{0.25,0.63,0.44}{#1}}
\newcommand{\BuiltInTok}[1]{\textcolor[rgb]{0.00,0.50,0.00}{#1}}
\newcommand{\CharTok}[1]{\textcolor[rgb]{0.25,0.44,0.63}{#1}}
\newcommand{\CommentTok}[1]{\textcolor[rgb]{0.38,0.63,0.69}{\textit{#1}}}
\newcommand{\CommentVarTok}[1]{\textcolor[rgb]{0.38,0.63,0.69}{\textbf{\textit{#1}}}}
\newcommand{\ConstantTok}[1]{\textcolor[rgb]{0.53,0.00,0.00}{#1}}
\newcommand{\ControlFlowTok}[1]{\textcolor[rgb]{0.00,0.44,0.13}{\textbf{#1}}}
\newcommand{\DataTypeTok}[1]{\textcolor[rgb]{0.56,0.13,0.00}{#1}}
\newcommand{\DecValTok}[1]{\textcolor[rgb]{0.25,0.63,0.44}{#1}}
\newcommand{\DocumentationTok}[1]{\textcolor[rgb]{0.73,0.13,0.13}{\textit{#1}}}
\newcommand{\ErrorTok}[1]{\textcolor[rgb]{1.00,0.00,0.00}{\textbf{#1}}}
\newcommand{\ExtensionTok}[1]{#1}
\newcommand{\FloatTok}[1]{\textcolor[rgb]{0.25,0.63,0.44}{#1}}
\newcommand{\FunctionTok}[1]{\textcolor[rgb]{0.02,0.16,0.49}{#1}}
\newcommand{\ImportTok}[1]{\textcolor[rgb]{0.00,0.50,0.00}{\textbf{#1}}}
\newcommand{\InformationTok}[1]{\textcolor[rgb]{0.38,0.63,0.69}{\textbf{\textit{#1}}}}
\newcommand{\KeywordTok}[1]{\textcolor[rgb]{0.00,0.44,0.13}{\textbf{#1}}}
\newcommand{\NormalTok}[1]{#1}
\newcommand{\OperatorTok}[1]{\textcolor[rgb]{0.40,0.40,0.40}{#1}}
\newcommand{\OtherTok}[1]{\textcolor[rgb]{0.00,0.44,0.13}{#1}}
\newcommand{\PreprocessorTok}[1]{\textcolor[rgb]{0.74,0.48,0.00}{#1}}
\newcommand{\RegionMarkerTok}[1]{#1}
\newcommand{\SpecialCharTok}[1]{\textcolor[rgb]{0.25,0.44,0.63}{#1}}
\newcommand{\SpecialStringTok}[1]{\textcolor[rgb]{0.73,0.40,0.53}{#1}}
\newcommand{\StringTok}[1]{\textcolor[rgb]{0.25,0.44,0.63}{#1}}
\newcommand{\VariableTok}[1]{\textcolor[rgb]{0.10,0.09,0.49}{#1}}
\newcommand{\VerbatimStringTok}[1]{\textcolor[rgb]{0.25,0.44,0.63}{#1}}
\newcommand{\WarningTok}[1]{\textcolor[rgb]{0.38,0.63,0.69}{\textbf{\textit{#1}}}}
\usepackage{longtable,booktabs,array}
\usepackage{calc} % for calculating minipage widths
% Correct order of tables after \paragraph or \subparagraph
\usepackage{etoolbox}
\makeatletter
\patchcmd\longtable{\par}{\if@noskipsec\mbox{}\fi\par}{}{}
\makeatother
% Allow footnotes in longtable head/foot
\IfFileExists{footnotehyper.sty}{\usepackage{footnotehyper}}{\usepackage{footnote}}
\makesavenoteenv{longtable}
\setlength{\emergencystretch}{3em} % prevent overfull lines
\providecommand{\tightlist}{%
  \setlength{\itemsep}{0pt}\setlength{\parskip}{0pt}}
\usepackage{bookmark}
\IfFileExists{xurl.sty}{\usepackage{xurl}}{} % add URL line breaks if available
\urlstyle{same}
\hypersetup{
  hidelinks,
  pdfcreator={LaTeX via pandoc}}

\author{}
\date{}

\begin{document}

{
\setcounter{tocdepth}{3}
\tableofcontents
}
\section{Trabajo Final del Curso}\label{trabajo-final-del-curso}

Curso: \texttt{Data\ Visualization}\\
Herramientas principales: \texttt{Tableau} y \texttt{Python} con
notebooks

\subsection{1. Propósito del trabajo
final}\label{1-propuxf3sito-del-trabajo-final}

El trabajo final del curso consiste en desarrollar un producto integral
de visualización de datos orientado a la toma de decisiones. El proyecto
debe mostrar que el equipo puede recorrer todo el flujo de trabajo visto
en clase: comprender un problema, evaluar una fuente de datos,
perfilarla, limpiarla, modelarla, analizarla, diseñar visualizaciones
adecuadas, construir un dashboard funcional y defender técnicamente sus
decisiones.

El objetivo no es solo presentar gráficos atractivos. El objetivo es
construir una solución analítica sólida, reproducible y bien
argumentada.

\subsection{2. Modalidad de trabajo}\label{2-modalidad-de-trabajo}

\begin{itemize}
\tightlist
\item
  Equipos de \texttt{2} a \texttt{3} estudiantes.
\item
  Si el docente lo aprueba, puede adaptarse a modalidad individual.
\item
  El equipo debe trabajar un solo tema durante todo el ciclo.
\item
  Cada entrega parcial debe reutilizar y mejorar la entrega anterior.
\end{itemize}

\subsection{3. Producto final esperado}\label{3-producto-final-esperado}

Al final del curso, cada equipo deberá entregar:

\begin{itemize}
\tightlist
\item
  Un \texttt{dashboard\ final} en \texttt{Tableau}.
\item
  Una \texttt{historia\ visual} breve o secuencia narrativa equivalente.
\item
  Un conjunto de \texttt{notebooks} en \texttt{Python} con el pipeline
  del proyecto.
\item
  Los archivos de datos preparados para Tableau.
\item
  Una bitácora técnica con las decisiones de preparación, validación y
  modelado.
\item
  Una presentación final para la defensa del proyecto.
\end{itemize}

\subsection{4. Qué debe resolver el
proyecto}\label{4-quuxe9-debe-resolver-el-proyecto}

El proyecto debe responder una pregunta analítica real y relevante. Esa
pregunta debe estar vinculada con una decisión, necesidad de monitoreo,
evaluación de desempeño, identificación de patrones o comprensión de un
fenómeno.

El equipo debe poder responder con claridad:

\begin{itemize}
\tightlist
\item
  ¿Cuál es el problema que se quiere analizar?
\item
  ¿Quién usará el dashboard?
\item
  ¿Qué decisión o acción podría apoyarse con este análisis?
\item
  ¿Por qué el dataset elegido permite responder esa pregunta?
\end{itemize}

\subsection{5. Requisitos del dataset}\label{5-requisitos-del-dataset}

El dataset debe cumplir, idealmente, con las siguientes condiciones:

\begin{itemize}
\tightlist
\item
  Tener al menos \texttt{2,000} registros.
\item
  Tener al menos \texttt{10} variables útiles.
\item
  Incluir al menos \texttt{1} variable temporal.
\item
  Incluir al menos \texttt{1} variable categórica relevante.
\item
  Incluir al menos \texttt{1} variable que permita segmentar.
\item
  Incluir al menos \texttt{4} variables numéricas.
\item
  De preferencia, incluir una dimensión geográfica.
\item
  Si se quiere aplicar \texttt{PCA} o \texttt{t-SNE}, contar con
  suficientes variables numéricas o una representación adecuada.
\end{itemize}

Además:

\begin{itemize}
\tightlist
\item
  Debe ser una fuente real, pública, institucional o verificable.
\item
  Debe tener trazabilidad: origen, periodo, cobertura y limitaciones.
\item
  No debe ser un dataset trivial.
\item
  No debe ser un dataset tan limpio que impida trabajar perfilado y
  limpieza.
\end{itemize}

\subsection{6. Cobertura mínima obligatoria del
curso}\label{6-cobertura-muxednima-obligatoria-del-curso}

El trabajo final debe evidenciar todos los siguientes componentes:

\begin{enumerate}
\def\labelenumi{\arabic{enumi}.}
\tightlist
\item
  Definición del problema y del usuario objetivo.
\item
  Perfilado del dato y comprensión de la granularidad.
\item
  Limpieza y preparación.
\item
  Modelado analítico o estructuración de la fuente.
\item
  Análisis exploratorio.
\item
  Selección y justificación de gráficos.
\item
  Segmentación e interpretación.
\item
  Construcción de métricas o cálculos analíticos.
\item
  Storytelling técnico.
\item
  Revisión de accesibilidad y diseño visual.
\item
  Visualización longitudinal.
\item
  Visualización transversal o comparativa.
\item
  Componente avanzado.
\item
  Construcción del dashboard final.
\item
  Defensa metodológica y visual.
\end{enumerate}

\subsection{7. Configuración mínima del producto
final}\label{7-configuraciuxf3n-muxednima-del-producto-final}

Como mínimo, el proyecto final debe contener:

\begin{itemize}
\tightlist
\item
  \texttt{1} problema analítico claramente formulado.
\item
  \texttt{1} usuario objetivo principal.
\item
  \texttt{1} dataset principal validado.
\item
  \texttt{1} pipeline reproducible en \texttt{Python}.
\item
  \texttt{1} workbook en \texttt{Tableau} con al menos \texttt{8} hojas
  analíticas útiles.
\item
  \texttt{1} dashboard final con al menos \texttt{3} bloques
  funcionales:

  \begin{itemize}
  \tightlist
  \item
    contexto y KPIs
  \item
    exploración comparativa o segmentada
  \item
    módulo temporal, comparativo o avanzado
  \end{itemize}
\item
  \texttt{1} historia visual de \texttt{3} a \texttt{5} pantallas o una
  secuencia equivalente.
\item
  \texttt{1} vista temporal sólida.
\item
  \texttt{1} vista transversal sólida.
\item
  \texttt{1} componente avanzado integrado o anexado con justificación
  metodológica.
\item
  \texttt{1} documento breve de QA con validación técnica, limitaciones
  y decisiones de diseño.
\end{itemize}

\subsection{8. Entregas parciales}\label{8-entregas-parciales}

El trabajo final se desarrollará de manera progresiva mediante entregas
parciales.

\subsubsection{Entrega 1. Propuesta del
proyecto}\label{entrega-1-propuesta-del-proyecto}

Semana sugerida: \texttt{2}

Debe incluir:

\begin{itemize}
\tightlist
\item
  tema del proyecto
\item
  pregunta analítica principal
\item
  usuario objetivo
\item
  fuente de datos propuesta
\item
  hipótesis iniciales
\item
  valor potencial del proyecto
\end{itemize}

Entregables:

\begin{itemize}
\tightlist
\item
  documento breve en \texttt{PDF}
\item
  enlace o archivo del dataset
\item
  ficha resumen del proyecto
\end{itemize}

Criterios mínimos:

\begin{itemize}
\tightlist
\item
  la pregunta analítica no es trivial
\item
  el usuario objetivo está claramente definido
\item
  el dataset permite temporalidad, segmentación y comparación
\item
  se explican riesgos iniciales de calidad o cobertura
\end{itemize}

\subsubsection{Entrega 2. Perfilado, diccionario y limpieza
inicial}\label{entrega-2-perfilado-diccionario-y-limpieza-inicial}

Semana sugerida: \texttt{4}

Debe incluir:

\begin{itemize}
\tightlist
\item
  perfilado del dataset
\item
  unidad de análisis
\item
  granularidad
\item
  diccionario de datos
\item
  principales problemas de calidad
\item
  reglas de limpieza
\item
  bitácora inicial
\end{itemize}

Entregables:

\begin{itemize}
\tightlist
\item
  notebook de perfilado y limpieza
\item
  tabla de perfilado
\item
  dataset limpio preliminar
\item
  bitácora de transformaciones
\end{itemize}

Criterios mínimos:

\begin{itemize}
\tightlist
\item
  se reportan nulos, cardinalidad, duplicados y campos críticos
\item
  se justifican las decisiones de limpieza
\item
  el dataset limpio puede conectarse a Tableau sin ambigüedad
\item
  se diferencia claramente la fuente original de la fuente transformada
\end{itemize}

\subsubsection{Entrega 3. Exploración, chart selection e insights
iniciales}\label{entrega-3-exploraciuxf3n-chart-selection-e-insights-iniciales}

Semana sugerida: \texttt{6}

Debe incluir:

\begin{itemize}
\tightlist
\item
  primeras visualizaciones exploratorias
\item
  comparación, distribución, relación y tendencia
\item
  criterios de selección de gráficos
\item
  primeros insights
\end{itemize}

Entregables:

\begin{itemize}
\tightlist
\item
  notebook exploratorio
\item
  workbook preliminar en Tableau
\item
  documento corto con \texttt{3} a \texttt{5} insights
\end{itemize}

Criterios mínimos:

\begin{itemize}
\tightlist
\item
  se muestran distintas familias de gráficos
\item
  cada gráfico tiene una justificación técnica breve
\item
  al menos \texttt{2} opciones visuales fueron descartadas y explicadas
\item
  los insights están redactados con lenguaje analítico
\end{itemize}

\subsubsection{Entrega 4. Modelado, métricas y dashboard
alpha}\label{entrega-4-modelado-muxe9tricas-y-dashboard-alpha}

Semana sugerida: \texttt{8}

Debe incluir:

\begin{itemize}
\tightlist
\item
  estructura analítica o relacional
\item
  validación de joins o relationships
\item
  métricas derivadas
\item
  segmentación
\item
  parámetros, filtros o lógica analítica
\item
  primera versión funcional del dashboard
\end{itemize}

Entregables:

\begin{itemize}
\tightlist
\item
  notebook de modelado y cálculos
\item
  fuentes analíticas para Tableau
\item
  dashboard alpha
\end{itemize}

Criterios mínimos:

\begin{itemize}
\tightlist
\item
  la estructura relacional está validada
\item
  las métricas responden a la pregunta del proyecto
\item
  el dashboard alpha ya tiene flujo de lectura
\item
  existe al menos una vista segmentada útil
\end{itemize}

\subsubsection{Entrega 5. Storytelling, accesibilidad y módulo
temporal/comparativo}\label{entrega-5-storytelling-accesibilidad-y-muxf3dulo-temporalcomparativo}

Semana sugerida: \texttt{10}

Debe incluir:

\begin{itemize}
\tightlist
\item
  rediseño narrativo del dashboard
\item
  mejora de accesibilidad
\item
  vista longitudinal sólida
\item
  vista transversal sólida
\end{itemize}

Entregables:

\begin{itemize}
\tightlist
\item
  dashboard revisado
\item
  historia visual breve
\item
  checklist de accesibilidad
\item
  reflexión metodológica corta
\end{itemize}

Criterios mínimos:

\begin{itemize}
\tightlist
\item
  mejora visible respecto a la versión alpha
\item
  títulos, color, contraste y jerarquía están mejor resueltos
\item
  la vista temporal y la vista transversal son defendibles
\item
  la historia visual comunica un mensaje, no solo muestra gráficos
\end{itemize}

\subsubsection{Entrega 6. Componente avanzado y dashboard
beta}\label{entrega-6-componente-avanzado-y-dashboard-beta}

Semana sugerida: \texttt{12}

Debe incluir:

\begin{itemize}
\tightlist
\item
  aplicación de \texttt{PCA}, \texttt{t-SNE} o técnica equivalente
  aprobada
\item
  explicación metodológica
\item
  integración del resultado al análisis
\item
  versión beta del dashboard
\end{itemize}

Entregables:

\begin{itemize}
\tightlist
\item
  notebook avanzado
\item
  exportables para Tableau
\item
  vista avanzada integrada o anexa
\item
  dashboard beta
\end{itemize}

Criterios mínimos:

\begin{itemize}
\tightlist
\item
  la técnica está correctamente explicada
\item
  se documentan variables, parámetros y limitaciones
\item
  el resultado aporta valor analítico al caso
\item
  el dashboard beta ya está listo para revisión final
\end{itemize}

\subsubsection{Entrega 7. Entrega final y
defensa}\label{entrega-7-entrega-final-y-defensa}

Semana sugerida: \texttt{14}

Debe incluir:

\begin{itemize}
\tightlist
\item
  dashboard final
\item
  historia visual o secuencia narrativa
\item
  QA técnico
\item
  defensa integral del pipeline
\end{itemize}

Entregables:

\begin{itemize}
\tightlist
\item
  dashboard final en Tableau
\item
  carpeta final de notebooks
\item
  base final preparada
\item
  presentación de defensa
\item
  resumen ejecutivo
\end{itemize}

Criterios mínimos:

\begin{itemize}
\tightlist
\item
  el dashboard responde claramente la pregunta principal
\item
  el pipeline técnico es reproducible
\item
  el equipo demuestra cobertura de todo lo visto en clase
\item
  la defensa muestra dominio de supuestos, límites y decisiones
\end{itemize}

\subsection{9. Cronograma resumido}\label{9-cronograma-resumido}

\begin{longtable}[]{@{}lll@{}}
\toprule\noalign{}
Semana & Tema del curso & Hito del proyecto \\
\midrule\noalign{}
\endhead
\bottomrule\noalign{}
\endlastfoot
1 & Introducción y ecosistema & Exploración de temas \\
2 & Perfilado y granularidad & Entrega \texttt{1} \\
3 & Limpieza y preparación & Avance técnico \\
4 & Modelado y fuentes & Entrega \texttt{2} \\
5 & Visualización exploratoria & Prototipos \\
6 & Segmentación e insights & Entrega \texttt{3} \\
7 & Cálculos analíticos & Construcción de métricas \\
8 & Storytelling técnico & Entrega \texttt{4} \\
9 & Accesibilidad y diseño & Revisión visual \\
10 & Series temporales & Entrega \texttt{5} \\
11 & Comparación transversal & Integración comparativa \\
12 & \texttt{PCA} y \texttt{t-SNE} & Entrega \texttt{6} \\
13 & Dashboard engineering y QA & Cierre técnico \\
14 & Capstone y defensa & Entrega \texttt{7} \\
\end{longtable}

\subsection{10. Criterios globales de
evaluación}\label{10-criterios-globales-de-evaluaciuxf3n}

La evaluación del trabajo final considerará, como mínimo, los siguientes
aspectos:

\begin{itemize}
\tightlist
\item
  definición del problema y claridad del objetivo
\item
  calidad, trazabilidad y perfilado del dato
\item
  limpieza y preparación reproducible
\item
  modelado y consistencia técnica
\item
  rigor analítico e insights
\item
  calidad visual y accesibilidad
\item
  storytelling y comunicación
\item
  componente avanzado
\item
  defensa final
\end{itemize}

\subsection{11. Formato de entrega
sugerido}\label{11-formato-de-entrega-sugerido}

Se recomienda organizar la entrega final de esta forma:

\begin{Shaded}
\begin{Highlighting}[]
\NormalTok{proyecto{-}final/}
\NormalTok{  data/}
\NormalTok{  notebooks/}
\NormalTok{  outputs/}
\NormalTok{  tableau/}
\NormalTok{  docs/}
\NormalTok{  README.md}
\end{Highlighting}
\end{Shaded}

Contenido esperado:

\begin{itemize}
\tightlist
\item
  \texttt{data/}: fuentes originales y fuentes limpias
\item
  \texttt{notebooks/}: notebooks de trabajo
\item
  \texttt{outputs/}: exportables para Tableau
\item
  \texttt{tableau/}: workbook o enlace publicado
\item
  \texttt{docs/}: diccionario, bitácora, QA, presentación y resumen
  ejecutivo
\item
  \texttt{README.md}: explicación general del proyecto
\end{itemize}

\subsection{12. Relación con los laboratorios del
curso}\label{12-relaciuxf3n-con-los-laboratorios-del-curso}

Cada entrega parcial debe apoyarse en los laboratorios y notebooks
semanales del curso. La idea es que el trabajo final no se construya al
final del semestre como un esfuerzo aislado, sino como una acumulación
ordenada de evidencia técnica y analítica.

Se recomienda usar como base:

\begin{itemize}
\tightlist
\item
  perfilado y granularidad
\item
  limpieza y preparación
\item
  modelado de fuentes
\item
  chart selection
\item
  segmentación e insights
\item
  cálculos analíticos
\item
  storytelling y anotación
\item
  accesibilidad y diseño
\item
  análisis temporal
\item
  comparación transversal
\item
  \texttt{PCA} o \texttt{t-SNE}
\item
  dashboard engineering
\item
  QA final
\end{itemize}

\subsection{13. Recomendaciones
importantes}\label{13-recomendaciones-importantes}

\begin{itemize}
\tightlist
\item
  Elijan una pregunta concreta, no un tema demasiado amplio.
\item
  No construyan el dashboard antes de entender la granularidad del dato.
\item
  No usen gráficos por variedad visual; úsenlos porque responden mejor a
  una pregunta.
\item
  No oculten problemas de calidad; documentarlos también es parte del
  trabajo.
\item
  No usen el componente avanzado como adorno técnico; debe aportar al
  análisis.
\item
  Cuiden la legibilidad, el contraste y el orden visual del dashboard.
\item
  Ensayen la defensa con foco en decisiones, no solo en resultados.
\end{itemize}

\subsection{14. Qué se espera de una buena entrega
final}\label{14-quuxe9-se-espera-de-una-buena-entrega-final}

Una buena entrega final:

\begin{itemize}
\tightlist
\item
  responde una pregunta real
\item
  usa datos con trazabilidad
\item
  documenta su pipeline
\item
  construye insights defendibles
\item
  comunica con claridad
\item
  y puede sostener técnicamente lo que afirma
\end{itemize}

\end{document}
